\originalchapter{8.334作业6：超出自旋波、长程畴壁与熔化}
\psource{Mehran Kardar，8.334 Spring 2014，PS6，截止2014年5月7日。原选做题3保留。}
\begin{exercise}[原题1(a)--(h)]\label{p:334-6-1}
单位$n$分量场$s$，$\beta H=\int\dd^dx\dd^dy\,K(|x-y|)s(x)\cdot s(y)$，核Fourier为$Kq^\omega/2$、$0<\omega<2$。(a)长距尾幂；(b)$s=(\pi,\sigma)$、$\sigma=\sqrt{1-\pi^2}$后有效作用；(c)二次Fourier与协方差；(d)消窄壳后$\avg\sigma$及单位长度恢复$\zeta$；(e)利用非解析核不被局部微扰修正求$K'$；(f)$T=1/K$流；(g)$d=\omega+\epsilon$的$T_c,\nu$；(h)小场$-h\cdot\int s$的流和$y_h$。
\end{exercise}
\begin{solution}
尺度$q^\omega$对应实空间$r^{-d-\omega}$尾；正$q^\omega$核的非接触部分有铁磁所需负号，并有对角常数让均匀态能0。不要把Fourier正系数误读为正的反铁磁实空间尾。

在北极坐标$\sigma=\sqrt{1-\pi^2}$，球面测度为$\prod_x\dd^{n-1}\pi/\sqrt{1-\pi^2}$。可写精确有效能为原核代入加$\tfrac12\sum_x\log(1-\pi_x^2)$。展开二次动能$H_0=K\int_q q^\omega\pi(q)\cdot\pi(-q)/2$，非局域四次项$K\int_q q^\omega|(\pi^2)(q)|^2/8$，另有测度顶点；这些顶点与对称约束一起保证不能任意产生Goldstone质量。二次协方差$\avg{\pi_i(q)\pi_j(q')}_0=\delta_{ij}(2\pi)^d\delta(q+q')/(Kq^\omega)$。

定义$I_\omega=\int_{\Lambda/b<q<\Lambda}\dd^dq/(2\pi)^dq^{-\omega}$，壳平均$\avg\sigma_>=1-(n-1)TI_\omega/2+\cdots=\zeta$。慢场$s'=s_</\zeta$，距离縮放给$K'=b^{d-\omega}\zeta^2K$，非解析动能系数无额外局部顶点修正（按题给假设）。窄壳$I_\omega=A_d\Lambda^{d-\omega}\dd\ell$，
\begin{equation}
\dot T=-(d-\omega)T+(n-1)A_d\Lambda^{d-\omega}T^2+\cdots.
\end{equation}
$n>1$、$\epsilon=d-\omega>0$小时$T_c=\epsilon/[(n-1)A_d\Lambda^\epsilon]+O(\epsilon^2)$，临界不稳定斜率$\epsilon$，$\nu=1/\epsilon+O(1)$。小场$h'=b^d\zeta h$，在固定点$y_h=d-\epsilon/2=(d+\omega)/2+O(\epsilon^2)$，对应$\eta=2-\omega$。短程模型有额外动能修正导致$n-2$，这里是$n-1$；不可把两公式混抄。$n=1$没有连续横向自旋波，不能除以$n-1$。
\end{solution}
\begin{exercise}[原题2(a)--(c)]\label{p:334-6-2}
XY在$d=2+\epsilon$的题给流$\dot T=-\epsilon T+4\pi^3y^2$、$\dot y=(2-\pi/T)y$。(a)有限温固定点；(b)最低非零阶本征值；(c)保持最低非平凡$\epsilon$阶估$d=3$超流$\nu,\alpha$。
\end{exercise}
\begin{solution}
非零$y$令$T_* =\pi/2$、$y_*^2=\epsilon/(8\pi^2)$。Jacobian为$\left(\begin{smallmatrix}-\epsilon&8\pi^3y_*\\4y_*/\pi&0\end{smallmatrix}\right)$，本征值$[-\epsilon\pm\sqrt{\epsilon^2+16\epsilon}]/2=\pm2\sqrt\epsilon+O(\epsilon)$。所以领先$\nu=1/(2\sqrt\epsilon)$；用$\alpha=2-d\nu$外推$\epsilon=1,d=3$给$\nu\simeq0.5,\alpha\simeq0.5$。

这是近二维截断流的粗外推，$\epsilon=1$没有小参数控制，不能与精确三维值等同；也不能先保留某个高阶本征值再丢其他同阶流修正，给出貌似更精确的数字。$\epsilon\to0$的幂次展开不还原二维BKT有限$\nu$，二维本来是本质奇点。
\end{solution}
\begin{exercise}[原题3(a)--(c)，选做]\label{p:334-6-3}
XY加$-\beta H_p=h_p\int\dd^2x\cos(p\theta)$，$p=1$磁场，$p=2,3,4,6$可来自离散晶格对称。(a)无涡旋sine--Gordon流；(b)含涡旋时证明$\dot h_p=[2-p^2/(4\pi K)]h_p$、$\dot K^{-1}=-\pi p^2h_p^2K^{-2}/2+4\pi^3y^2$、$\dot y=(2-\pi K)y$；(c)解释$8\pi/p^2<K^{-1}<\pi/2$区间、$p>4,p<4$相图与$p=4$边缘情况。
\end{exercise}
\begin{solution}
Gaussian壳角方差$\avg{\theta_>^2}=\dd\ell/(2\pi K)$，顶点平均$\avg{\cos[p(\theta_<+\theta_>)]}=\e^{-p^2\dd\ell/(4\pi K)}\cos p\theta_<$，体积因子2给场流。二阶连接累积量中两个近邻余弦取差角，$\cos[p(\theta(x)-\theta(y))]\simeq1-p^2[(x-y)\cdot\nabla\theta]^2/2$；角平均与截断壳积分按题给归一化给$\dot K=\pi p^2h_p^2/2$，即角被锁定使刚度增大。涡旋偶极则屏蔽刚度，贡献$\dot K^{-1}=4\pi^3y^2$，将两项合并即(b)。$h_p$归一化改变二阶非普适系数，不改顶点尺度维。

两种扰动都无关需$2-p^2/(4\pi K)<0$且$2-\pi K<0$，所以$8\pi/p^2<K^{-1}<\pi/2$，只有$p>4$非空。这是Gaussian\textbf{稳定区}，不是微扰方程只能在该区推导：区外小初始耦合仍能用流判定其相关，长到$O(1)$后才失控。原题“only valid”须以此条件解释。

$p>4$典型升温为离散锁定序、代数关联、涡旋无序两次转变；$p=2,3$无稳定中间Gaussian相，一般一次离散序转变；$p=1$外场去掉所有对称破缺，不保证有真转变。$p=4$二扰动可边缘竞争，关联到Ashkin--Teller类连续变指数，截断线性流不足以全定相图。
\begin{center}\begin{tikzpicture}[x=1.1cm,y=.9cm]\draw[->](0,0)--(9,0)node[right]{$T$};\node[left]at(0,0){$p>4$};\draw[thick](3,-.15)--(3,.15);\draw[thick](6,-.15)--(6,.15);\node[above]at(1.4,0){离散序};\node[above]at(4.5,0){代数序};\node[above]at(7.5,0){无序};\draw[->](0,-1.4)--(9,-1.4)node[right]{$T$};\node[left]at(0,-1.4){$p=2,3$};\draw[thick](4.5,-1.55)--(4.5,-1.25);\node[above]at(2.3,-1.4){离散序};\node[above]at(6.7,-1.4){无序};\end{tikzpicture}\end{center}
\end{solution}
\begin{exercise}[原题4(a)--(f)]\label{p:334-6-4}
一维标量$\pm1$自旋有逆平方铁磁相互作用，物理能$\beta H=-K\int\dd x\dd y\,s(x)s(y)/(2|x-y|^2)+\int\Phi(s)\dd x$，$\Phi$强制$\pm1$，例如$g(s^2-1)^2,g\gg1$。原卷把$\Phi$加在$-\beta H$中会反向失稳，本题取其所述“偏好$\pm1$”的稳定能量符号。(a)尺度$a$畴壁能；(b)单壁概率给$K_c$下界；(c)交替电荷畴壁对数气；(d)由$K_{\rm eff}=K-2Ky_0^2\int_a^\infty(r\dd r/a^2)(a/r)^{4K}+O(y_0^4)$得壳流；(e)画流及相；(f)推导此屏蔽式。原卷积分未显写$a^{-2}$，等价于所有长度按$a$计的约定。
\end{exercise}
\begin{solution}
单壁两侧相反的能代价$2K\int_0^L\dd u\int_0^L\dd v/(u+v)^2\simeq2K\log(L/a)$；位置熵$\log(L/a)$，单壁概率尺度$(L/a)^{1-2K}$，所以$K<1/2$壁大量出现，不可能保持该铁磁有序，$K_c\ge1/2$。

利用$1/(x-y)^2=\partial_x\partial_y\log|x-y|$并分部积分，$s'(x)=2\sum_iq_i\delta(x-x_i)$，得到$-\beta H_Q=4K\sum_{i<j}q_iq_j\log(|x_i-x_j|/a)+\log y_0\sum_i|q_i|$；电荷必须正负交替，局域壁核贡献逸度。

缩放$a\to ba$，中性$\sum_{i<j}q_iq_j=-N_w/2$给每壁$b^{-2K}$，位置积分$b$给$\dot y_0=(1-2K)y_0$。积分掉$a<r<ba$小壁对给$\dot K=-2Ky_0^2$。置$T=1/K$，$\dot T=2Ty_0^2$；$T<2$且小$y$时流向$y=0$有序固定线，$T>2$时壁相关流向高温无序，分界靠近$(T,y)=(2,0)$，具有类BKT本质尺度，强耦合处截断流不可靠。
\begin{center}\begin{tikzpicture}[x=2cm,y=2cm]\draw[->](0,0)--(2.5,0)node[right]{$T$};\draw[->](0,0)--(0,1.25)node[above]{$y_0$};\node[below]at(1.5,0){2};\draw[dashed](1.5,0)--(1.5,1);\draw[->,thick](.5,.65)..controls(.9,.2)..(1,.04);\draw[->,thick](1.3,.6)..controls(1.7,.75)..(2.2,1.1);\draw[->,thick](1.9,.25)..controls(2,.5)..(2.2,.7);\node at(.55,1){有序};\node at(1.9,1.25){无序};\end{tikzpicture}\end{center}

(f)固定外壁$+,-$于$0,L$，区间内允许小对方向$-,+$于$x,x+r$，相对权重$(y_0^2/a^2)(a/r)^{4K}$。与外壁交互的指数增量为$4K\log[(x+r)/x]+4K\log[(L-x)/(L-x-r)]\simeq4Kr[1/x+1/(L-x)]$。对$x$从核截断$a$积分到$L-a$给$8Kr\log(L/a)$；因此$\Delta\log W=8Ky_0^2\log(L/a)\int(r\dd r/a^2)(a/r)^{4K}$，与原$\log W=-4K\log(L/a)$比较正得题式。体积常数和端点核项吸收于自由能与逸度；交替约束允许一阶$r$贡献，是它比二维任意方向偶极屏蔽更简单的原因。
\end{solution}
\begin{exercise}[原题5(a)--(d)]\label{p:334-6-5}
各向同性弹性$\beta H=\tfrac12\int[2\mu u_{ij}u_{ij}+\lambda u_{ii}u_{jj}]\dd^dx$，$u_{ij}=(\partial_iu_j+\partial_ju_i)/2$。（原题将正稳定弹性能写在$-\beta H$之后，改正此符号。）(a)Fourier模；(b)协方差；(c)晶格截断$a$下$U^2(x)=\avg{|u(x)-u(0)|^2}$；(d)$\mu=\widehat\mu/(\kb T),\lambda=\widehat\lambda/(\kb T)$，由远距$U=c_La$的Lindemann准则求$T_m$及维数依赖。
\end{exercise}
\begin{solution}
纵横投影$P^L_{ij}=q_iq_j/q^2$、$P^T=I-P^L$，弹性核$q^2[\mu P^T+(\lambda+2\mu)P^L]$。静态本征值横$\mu q^2$（$d-1$支）、纵$(\lambda+2\mu)q^2$；若加质量密度动能，声速$\sqrt{\widehat\mu/\rho},\sqrt{(\widehat\lambda+2\widehat\mu)/\rho}$。稳定需$\mu>0$、体积模量$\lambda+2\mu/d>0$。

协方差为$(2\pi)^d\delta(q+q')q^{-2}[P^T/\mu+P^L/(\lambda+2\mu)]$。故
\begin{equation}
U^2(x)=2\int_q\frac{1-\cos(qx)}{q^2}\left[\frac{d-1}{\mu}+\frac1{\lambda+2\mu}\right].
\end{equation}
$d>2$远距为$2A_d\Lambda^{d-2}\kb T[(d-1)/\widehat\mu+1/(\widehat\lambda+2\widehat\mu)]/(d-2)$，$\Lambda\sim a^{-1}$。所以
\begin{equation}
T_m=\frac{c_L^2a^2(d-2)}{2\kb A_d\Lambda^{d-2}[(d-1)/\widehat\mu+1/(\widehat\lambda+2\widehat\mu)]}.
\end{equation}
这是截断与经验$c_L$决定的估计，不是精确转变温。$d=2$有$\log r$、$d<2$有$r^{2-d}$发散，此\textbf{绝对远距位移}准则会给所有正温$T_m=0$。实际二维可有代数晶序和有限温缺陷熔化，需第\ref{ch:melting}章的相对序与位错方法；不能以此准则错误宣称所有二维材料无固相。
\end{solution}
